%%%%%%%%%%%%%%%%%%%%%%%%%%%%%%%%%%%%%%%%%%%%%%%%%%%%%
% YahooAnswers-XML — Entrega Final
% Bases de Datos Documentales — 25 de marzo 2026
% Estilo: dos columnas, IEEE-like
%%%%%%%%%%%%%%%%%%%%%%%%%%%%%%%%%%%%%%%%%%%%%%%%%%%%%
\documentclass[10pt,twocolumn,a4paper]{article}

% ── Layout IEEE-like ──
\usepackage[a4paper,top=1.9cm,bottom=2.4cm,left=1.6cm,right=1.6cm,columnsep=0.65cm]{geometry}

% ── Encoding ──
\usepackage[utf8]{inputenc}
\usepackage[T1]{fontenc}
\usepackage[nil]{babel}

% ── Math / Symbols ──
\usepackage{amsmath,amssymb,amsfonts}
\usepackage{textcomp}

% ── Graphics / Color ──
\usepackage{graphicx}
\usepackage{xcolor}
\definecolor{statusgreen}{rgb}{0.13,0.55,0.13}
\definecolor{codegray}{rgb}{0.5,0.5,0.5}

% ── Tables ──
\usepackage{booktabs}
\usepackage{array}
\usepackage{tabularx}
\usepackage{float}
\usepackage{caption}
\captionsetup{font=footnotesize,labelfont=bf,skip=3pt}

% ── URLs ──
\usepackage{url}
\usepackage[hidelinks]{hyperref}

% ── Lists ──
\usepackage{enumitem}

% ── Code listings ──
\usepackage{listings}
\lstset{
    basicstyle=\ttfamily\scriptsize,
    breaklines=true,
    frame=single,
    numbers=left,
    numberstyle=\tiny\color{codegray},
    keywordstyle=\color{blue!70!black},
    commentstyle=\color{green!50!black},
    stringstyle=\color{red!60!black},
    showstringspaces=false,
    tabsize=2,
    captionpos=b,
    literate={a'}{a}1 {e'}{e}1 {i'}{i}1 {o'}{o}1 {u'}{u}1
             {ny}{{n}}1
}
\lstdefinelanguage{XQuery}{
    morekeywords={let,for,where,return,if,then,else,and,or,not,
                  collection,doc,distinct-values,count,avg,min,max,
                  sum,string-length,contains,lower-case,subsequence,
                  format-number,current-dateTime,xs,string,number,
                  order,by,descending,ascending,in},
    morecomment=[s]{(:}{:)},
    morestring=[b]",
    sensitive=true
}
\lstdefinelanguage{XMLL}{
    morestring=[b]",
    morecomment=[s]{<!--}{-->},
    stringstyle=\color{red!60!black},
    keywordstyle=\color{blue!70!black}
}

% ── Section headings IEEE style ──
\usepackage{titlesec}
\titleformat{\section}{\normalsize\bfseries\centering}{\Roman{section}.}{0.5em}{\MakeUppercase}
\titleformat{\subsection}{\small\bfseries\itshape}{\Alph{subsection}.}{0.5em}{}
\titleformat{\subsubsection}{\small\itshape}{\arabic{subsubsection})}{0.4em}{}
\titlespacing{\section}{0pt}{8pt}{4pt}
\titlespacing{\subsection}{0pt}{5pt}{2pt}
\titlespacing{\subsubsection}{0pt}{3pt}{1pt}

% ── Abstract ──
\usepackage{abstract}
\renewcommand{\abstractnamefont}{\normalfont\small\bfseries}
\renewcommand{\abstracttextfont}{\normalfont\small\itshape}
\setlength{\absleftindent}{0pt}
\setlength{\absrightindent}{0pt}

% ── Header / Footer ──
\usepackage{fancyhdr}
\pagestyle{fancy}
\fancyhf{}
\renewcommand{\headrulewidth}{0pt}
\fancyfoot[C]{\small\thepage}

% ── IEEE-like author helpers ──
\newcommand{\IEEEauthorblockN}[1]{\textbf{#1}\\[2pt]}
\newcommand{\IEEEauthorblockA}[1]{{\footnotesize #1}}

% ── Paragraph spacing ──
\setlength{\parskip}{2pt}
\setlength{\parindent}{1em}

%%%%%%%%%%%%%%%%%%%%%%%%%%%%%%%%%%%%%%%%%%%%%%%%%%%%%
\begin{document}
%%%%%%%%%%%%%%%%%%%%%%%%%%%%%%%%%%%%%%%%%%%%%%%%%%%%%

% ── Title (one-column) ──
\onecolumn
\begin{center}
  {\large\bfseries YahooAnswers-XML: An\'alisis de Yahoo! Answers\\[4pt]
   con XML, XSD, XSLT, XQuery y eXist-db}\\[10pt]
  \begin{tabular}{cc}
    \textbf{Matteo Amagliani} & \textbf{Manuel Avil\'es Rodr\'iguez} \\[2pt]
    {\small Universidad de Castilla-La Mancha} & {\small Universidad de Castilla-La Mancha} \\
    {\small Grado en Ingenier\'ia Inform\'atica} & {\small Grado en Ingenier\'ia Inform\'atica} \\
    {\small matteo.amagliani1@alu.uclm.es} & {\small manuel.aviles@alu.uclm.es}
  \end{tabular}\\[10pt]
  \rule{0.93\textwidth}{0.4pt}
\end{center}

\begin{abstract}
\noindent
Este documento presenta el dise\~no e implementaci\'on de \textbf{YahooAnswers-XML},
un sistema de an\'alisis del corpus p\'ublico de Yahoo! Answers (octubre~2007)
mediante tecnolog\'ias XML y bases de datos documentales. El sistema procesa
16.952 preguntas y 161.365 respuestas distribuidas en cuatro categor\'ias
tem\'aticas bilingues (espa\~nol e italiano), aplicando un flujo completo
de extracci\'on, validaci\'on, transformaci\'on y consulta. Se describen el
dise\~no del esquema XSD, las transformaciones XSLT, la carga en
\textbf{eXist-db} y las cuatro consultas XQuery desarrolladas para el
an\'alisis de Redes Sociales y Calidad de Datos. Los resultados revelan
patrones de baja simetr\'ia usuario-pregunta/respuesta, diferencias
significativas entre comunidades hispanohablante e italohablante, y una
calidad de datos generalmente alta con algunas anomal\'ias identificadas.
\end{abstract}

\smallskip
\noindent{\footnotesize\textbf{Palabras clave---}\itshape
XML, XSD, XSLT, XQuery, eXist-db, Bases de Datos Documentales,
Yahoo! Answers, An\'alisis de Redes Sociales, Calidad de Datos.}
\vspace{6pt}

\begin{center}\rule{0.93\textwidth}{0.4pt}\end{center}
\vspace{4pt}
\twocolumn

%% ============================================
\section{Introducci\'on}
%% ============================================

Yahoo! Answers fue durante a\~nos uno de los sistemas de preguntas y
respuestas comunitarias m\'as grandes del mundo. El corpus p\'ublico
\textit{ydata-yanswers-all-questions-v1\_0}~\cite{yahoowebscope}, con
\textbf{4.483.032 preguntas} recogidas en octubre de 2007, contiene el
texto de cada pregunta, todas sus respuestas y la identificaci\'on de
la mejor respuesta seleccionada. Este proyecto aplica las tecnolog\'ias
del m\'odulo XML de Bases de Datos Documentales a cuatro categor\'ias
bilingues de dicho corpus.

\subsection{Objetivo general}

Analizar el corpus en las categor\'ias \textit{Pol\'itica y Gobierno}
y \textit{Noticias y Actualidad} (ES/IT) mediante XML y eXist-db,
realizando an\'alisis de Redes Sociales y Calidad de Datos.

\subsection{Objetivos espec\'ificos}

\begin{enumerate}[leftmargin=1.8em,itemsep=0pt,topsep=2pt]
    \item Dise\~nar un esquema XSD que valide la estructura del corpus.
    \item Construir un XSD ampliado con m\'etricas de calidad y
          anotaciones de red social.
    \item Implementar transformaciones XSLT a HTML navegable.
    \item Cargar los XML en eXist-db y desarrollar consultas XQuery.
    \item Analizar Redes Sociales y Calidad de Datos sobre el corpus.
\end{enumerate}

%% ============================================
\section{Corpus y Fuentes de Datos}
%% ============================================

\subsection{Yahoo! Answers Corpus (oct. 2007)}

El dataset utilizado es el corpus oficial de Yahoo! Research
Webscope~\cite{yahoowebscope}, un fichero XML de 12~GB. Para este trabajo
se extrajeron subconjuntos por categor\'ia mediante un \textit{parser} SAX
en Python, obteniendo cuatro ficheros XML manejables (Tabla~\ref{tab:categorias}).

\begin{table}[H]
\centering
\caption{Categor\'ias seleccionadas y volumen de datos.}
\label{tab:categorias}
\begin{tabular}{p{2.2cm}lrr}
\toprule
\textbf{Categor\'ia} & \textbf{Id.} & \textbf{Preg.} & \textbf{Resp.} \\
\midrule
Pol\'itica y Gob. & ES & 12.976 & 119.262 \\
Noticias y Act.   & ES &    674 &   4.740 \\
Politica e gov.   & IT &  2.118 &  25.456 \\
Notizie ed ev.    & IT &  1.184 &  11.907 \\
\midrule
\textbf{Total} & & \textbf{16.952} & \textbf{161.365} \\
\bottomrule
\end{tabular}
\end{table}

\subsection{Estructura XML del corpus}

El Listado~\ref{lst:xml} muestra la estructura de un documento.
Los campos \texttt{content} y \texttt{vot\_date} son opcionales:
ausentes en el 31,3\% y el 53,2\% de los documentos respectivamente.

\begin{lstlisting}[language=XMLL,caption={Estructura de un documento del corpus.},label=lst:xml]
<preguntas categoria="Politica y Gobierno">
 <document>
  <uri>3086062</uri>
  <subject>Pregunta del usuario</subject>
  <content>Detalle (opcional)</content>
  <bestanswer>Mejor respuesta</bestanswer>
  <nbestanswers>
   <answer_item>Resp. 1</answer_item>
  </nbestanswers>
  <maincat>Politica y Gobierno</maincat>
  <date>1163158230</date>
  <id>u774726</id>
  <best_id>u1011708</best_id>
 </document>
</preguntas>
\end{lstlisting}

%% ============================================
\section{Herramientas y Tecnolog\'ias}
%% ============================================

\begin{table}[H]
\centering
\caption{Tecnolog\'ias utilizadas.}
\label{tab:tech}
\begin{tabular}{ll}
\toprule
\textbf{Componente} & \textbf{Tecnolog\'ia} \\
\midrule
BD documental & eXist-db 6.x \\
Consultas & XQuery 3.1 \\
Esquema & XSD 1.1 \\
Transformaci\'on & XSLT 1.0 \\
Validaci\'on local & Python / lxml \\
Extracci\'on corpus & Python SAX \\
Control versiones & Git / GitHub \\
\bottomrule
\end{tabular}
\end{table}

\subsection{Estructura del repositorio}

\begin{lstlisting}[caption={Estructura del repositorio.}]
YahooAnswers-XML/
|-- src/
|   |-- extraer_categorias.py
|   +-- extraer_preguntas.py
|-- existdb/
|   |-- cargar_existdb.py
|   +-- ejecutar_xquery.py
|-- xquery/   (4 consultas .xq)
|-- xsd/      yahoo_answers.xsd
|-- xslt/     es_html.xsl, it_html.xsl
+-- data/     preguntas_*.xml
\end{lstlisting}

%% ============================================
\section{Dise\~no del Esquema XML (XSD)}
%% ============================================

\subsection{XSD original}

El fichero \texttt{yahoo\_answers.xsd} describe formalmente la estructura
del corpus. La Tabla~\ref{tab:xsd_campos} detalla los campos principales.

\begin{table}[H]
\centering
\caption{Campos del tipo \texttt{DocumentType}.}
\label{tab:xsd_campos}
\begin{tabular}{p{1.9cm}p{1.4cm}p{0.6cm}p{1.0cm}}
\toprule
\textbf{Campo} & \textbf{Tipo XSD} & \textbf{Ob.} & \textbf{Nota} \\
\midrule
\texttt{uri}         & posInt.  & S\'i & ID \'unico \\
\texttt{subject}     & string   & S\'i & Pregunta \\
\texttt{content}     & string   & No  & Detalle \\
\texttt{bestanswer}  & string   & S\'i & Mejor resp. \\
\texttt{nbestanswers}& ListaR.  & S\'i & Todas resp. \\
\texttt{date}        & long     & S\'i & Timestamp \\
\texttt{vot\_date}   & long     & No  & Fecha voto \\
\texttt{id}          & UserId   & S\'i & Autor \\
\texttt{best\_id}    & UserId   & S\'i & Respond. \\
\bottomrule
\end{tabular}
\end{table}

Se definen dos tipos simples restringidos: \texttt{CodigoIdioma}
(patr\'on \texttt{[a-z]\{2\}}) y \texttt{UserId} (patr\'on
\texttt{u[0-9]+}).

\begin{lstlisting}[language=XMLL,caption={Fragmento del XSD original.},label=lst:xsd]
<xs:complexType name="DocumentType">
 <xs:all>
  <xs:element name="uri"
   type="xs:positiveInteger"/>
  <xs:element name="content"
   type="xs:string" minOccurs="0"/>
  <xs:element name="id" type="UserId"/>
 </xs:all>
</xs:complexType>
<xs:simpleType name="UserId">
 <xs:restriction base="xs:string">
  <xs:pattern value="u[0-9]+"/>
 </xs:restriction>
</xs:simpleType>
\end{lstlisting}

\subsection{XSD ampliado}

Para el an\'alisis se dise\~n\'o un esquema extendido con dos nuevos
elementos opcionales:

\begin{itemize}[leftmargin=1.5em,itemsep=0pt,topsep=2pt]
    \item \textbf{\texttt{<metricas\_calidad>}}: campos derivados como
          \texttt{longitud\_bestanswer}, \texttt{num\_respuestas} y
          \texttt{tiempo\_resolucion\_seg}.
    \item \textbf{\texttt{<red\_social>}}: anotaciones para grafos sociales
          con \texttt{autor\_es\_respondedor}, \texttt{subcats\_usuario}
          y \texttt{peso\_enlace}.
\end{itemize}

Todos usan \texttt{minOccurs="0"} para mantener compatibilidad.

%% ============================================
\section{Transformaciones XSLT}
%% ============================================

Se dise\~naron dos transformaciones XSLT~1.0 (\texttt{es\_html.xsl}
e \texttt{it\_html.xsl}) que convierten los XML en p\'aginas HTML
navegables. Cada plantilla genera: cabecera con nombre de categor\'ia
y total de preguntas (\texttt{count(//document)}), una tarjeta por
documento con el \texttt{subject} como t\'itulo, el \texttt{bestanswer}
resaltado en verde, y un elemento \texttt{<details>} desplegable para
las dem\'as respuestas.

\begin{lstlisting}[language=XMLL,caption={Plantilla XSLT por documento.},label=lst:xslt]
<xsl:template match="document">
 <div class="documento">
  <div class="pregunta">
   <xsl:value-of select="subject"/>
  </div>
  <xsl:if test="content != ''">
   <div class="detalle">
    <xsl:value-of select="content"/>
   </div>
  </xsl:if>
  <div class="mejor-respuesta">
   <xsl:value-of select="bestanswer"/>
  </div>
 </div>
</xsl:template>
\end{lstlisting}

Se us\'o XSLT~1.0 para garantizar compatibilidad con \texttt{lxml} sin necesidad de Saxon.

%% ============================================
\section{Carga en eXist-db}
%% ============================================

Los cuatro ficheros XML se organizan en dos subcolecciones:

\begin{lstlisting}[caption={Colecciones en eXist-db.}]
/db/yahooanswers/
    es/  politica_y_gobierno.xml (95MB)
         noticias_y_actualidad.xml
    it/  politica_e_governo.xml
         notizie_ed_eventi.xml
\end{lstlisting}

El script \texttt{cargar\_existdb.py} usa HTTP~PUT sobre la REST~API.
Se aument\'o el l\'imite del servidor Jetty a 200~MB para acomodar
el fichero de mayor tama\~no. La verificaci\'on mediante
\texttt{count(collection("/db/yahooanswers")//document)}
devuelve 16.952, coincidiendo con el total esperado.

%% ============================================
\section{Consultas XQuery}
%% ============================================

\subsection{Estad\'isticas generales}

\begin{lstlisting}[language=XQuery,caption={Estadisticas por categoria.}]
let $es := collection(
  "/db/yahooanswers/es")//document
let $it := collection(
  "/db/yahooanswers/it")//document
return <estadisticas>{
  for $cat in distinct-values($es/maincat)
  let $docs := $es[maincat = $cat]
  order by count($docs) descending
  return <categoria nombre="{$cat}">
    <total>{count($docs)}</total>
    <media>{format-number(
      count($docs//answer_item)
      div count($docs),"0.00")}
    </media>
  </categoria>
}</estadisticas>
\end{lstlisting}

\begin{table}[H]
\centering
\caption{Estad\'isticas por categor\'ia (resultados reales).}
\label{tab:stats}
\begin{tabular}{p{2.0cm}rrr}
\toprule
\textbf{Categor\'ia} & \textbf{Docs} & \textbf{Resp.} & \textbf{Med.} \\
\midrule
Pol./Gob. (ES) & 12.976 & 119.262 & 9,19 \\
Noticias (ES)  &    674 &   4.740 & 7,03 \\
Politica (IT)  &  2.118 &  25.456 & 12,02 \\
Notizie (IT)   &  1.184 &  11.907 & 10,06 \\
\bottomrule
\end{tabular}
\end{table}

\subsection{An\'alisis de Redes Sociales}

Esta consulta calcula la simetr\'ia de participaci\'on, los usuarios
\textit{top} y el \textit{clustering} por categor\'ia.

\begin{lstlisting}[language=XQuery,caption={Calculo de simetria.}]
let $todos := ($es, $it)
let $ids := distinct-values($todos/id)
let $best := distinct-values(
              $todos/best_id)
let $solo_preg :=
  $ids[not(. = $best)]
let $solo_resp :=
  $best[not(. = $ids)]
let $ambas := $ids[. = $best]
return <simetria>
  <total>{count(distinct-values(
    ($ids,$best)))}</total>
  <solo_preg total="{
    count($solo_preg)}"/>
  <solo_resp total="{
    count($solo_resp)}"/>
  <ambas total="{count($ambas)}"/>
</simetria>
\end{lstlisting}

Los resultados de la Tabla~\ref{tab:simetria} confirman la baja
simetr\'ia caracter\'istica de estos sistemas~\cite{shen2015}:
solo el 15,4\% de los usuarios participa en ambos roles.

\begin{table}[H]
\centering
\caption{Simetr\'ia global de participaci\'on (12.881 usuarios).}
\label{tab:simetria}
\begin{tabular}{lrr}
\toprule
\textbf{Rol} & \textbf{Usuarios} & \textbf{\%} \\
\midrule
Solo preguntan & 6.165 & 47,9 \\
Solo responden & 4.732 & 36,7 \\
Hacen ambas    & 1.984 & 15,4 \\
\midrule
\textbf{Total} & \textbf{12.881} & \textbf{100} \\
\bottomrule
\end{tabular}
\end{table}

La Tabla~\ref{tab:simetria_cat} muestra la simetr\'ia por categor\'ia.
Las comunidades italianas presentan mayor simetr\'ia. Solo 2 usuarios
participan activamente en ES e IT, confirmando que ambas comunidades
son \textbf{pr\'acticamente independientes}.

\begin{table}[H]
\centering
\caption{Simetr\'ia por categor\'ia.}
\label{tab:simetria_cat}
\begin{tabular}{p{2.1cm}rrr}
\toprule
\textbf{Categor\'ia} & \textbf{Preg.} & \textbf{Resp.} & \textbf{Sim.} \\
\midrule
Noticias (ES)  &   444 &   489 &  5,5\% \\
Notizie (IT)   &   826 &   859 & 10,2\% \\
Politica (IT)  & 1.022 &   994 & 15,6\% \\
Pol./Gob. (ES) & 6.141 & 4.919 & 14,8\% \\
\bottomrule
\end{tabular}
\end{table}

El usuario m\'as activo (\texttt{u1012406}) acumula 398 mejores
respuestas concentradas en 2 categor\'ias, confirmando el elevado
\textit{clustering} de los respondedores \textit{top}~\cite{shen2015}.

\subsection{Calidad de Datos}

\begin{lstlisting}[language=XQuery,caption={Analisis de completitud.}]
let $todos := ($es,$it)
return <completitud>{
  for $c in ("content","vot_date")
  let $aus := count($todos) -
    count($todos[*[name()=$c]!=''])
  return <campo nombre="{$c}"
    pct="{format-number($aus div
      count($todos)*100,'0.0')}%"/>
}</completitud>
\end{lstlisting}

\begin{table}[H]
\centering
\caption{Completitud de campos opcionales.}
\label{tab:completitud}
\begin{tabular}{lrr}
\toprule
\textbf{Campo} & \textbf{Ausentes} & \textbf{\%} \\
\midrule
\texttt{content}   & 5.309 & 31,3\% \\
\texttt{vot\_date} & 9.010 & 53,2\% \\
\bottomrule
\end{tabular}
\end{table}

\begin{table}[H]
\centering
\caption{Indicadores de calidad por categor\'ia.}
\label{tab:calidad_cat}
\begin{tabular}{p{1.9cm}rrr}
\toprule
\textbf{Cat.} & \textbf{Sin cont.} & \textbf{Long. B.} & \textbf{M.R.} \\
\midrule
Noticias (ES) & 32,5\% & 344,9 & 7,03 \\
Notizie  (IT) & 34,2\% & 338,7 & 10,06 \\
Politica (IT) & 24,2\% & 502,5 & 12,02 \\
Pol./Gob.(ES) & 32,2\% & 442,1 &  9,19 \\
\bottomrule
\end{tabular}
\end{table}

Las principales anomal\'ias detectadas son: \textbf{677 documentos}
con \texttt{bestanswer} $<$20 caracteres; \textbf{847 documentos}
(5,0\%) con una sola respuesta; \textbf{0 inconsistencias de fechas}
(\texttt{date}$>$\texttt{res\_date}); y el 16,8\% de preguntas
resueltas el mismo d\'ia de publicaci\'on.

\subsection{Consultas Avanzadas}

\subsubsection{Comparativa ES vs.\ IT}

\begin{table}[H]
\centering
\caption{Comparativa comunidades ES vs.\ IT.}
\label{tab:comparativa}
\begin{tabular}{p{2.4cm}ll}
\toprule
\textbf{M\'etrica} & \textbf{ES} & \textbf{IT} \\
\midrule
Media resp./preg.  & 9,08  & 11,32 \\
Long. bestanswer   & 437,3 & 443,8 \\
Sin \texttt{content} & 32,2\% & 27,8\% \\
Usuarios preg.     & 6.473 & 1.687 \\
Usuarios resp.     & 5.170 & 1.548 \\
Resueltas mismo d\'ia & 16,0\% & 19,9\% \\
\bottomrule
\end{tabular}
\end{table}

La comunidad italiana es m\'as activa (11,32 resp./preg. frente a 9,08)
y sus respuestas son ligeramente m\'as largas (443,8 vs 437,3 caracteres),
sugiriendo mayor elaboraci\'on en el contexto pol\'itico italiano.

\subsubsection{Palabras clave en pol\'itica}

\begin{table}[H]
\centering
\caption{T\'erminos frecuentes en preguntas de pol\'itica.}
\label{tab:keywords}
\begin{tabular}{llr}
\toprule
\textbf{T\'ermino} & \textbf{Id.} & \textbf{Menc.} \\
\midrule
presidente    & ES & 732 \\
gobierno      & ES & 429 \\
ley           & ES & 273 \\
governo       & IT & 112 \\
legge         & IT &  48 \\
inmigrante    & ES &  52 \\
inmigraci\'on & ES &  26 \\
immigrazione  & IT &   4 \\
\bottomrule
\end{tabular}
\end{table}

Las menciones de inmigraci\'on en espa\~nol (78) superan ampliamente
las italianas (6), reflejando mayor relevancia del tema en el contexto
hispanohablante. La consulta confirma adem\'as que \textbf{0 preguntas}
del corpus carecen de al menos una respuesta.

%% ============================================
\section{Validaci\'on}
%% ============================================

Los cuatro XML superan la validaci\'on contra \texttt{yahoo\_answers.xsd}
con \texttt{lxml} sin errores. Las decisiones clave:
\begin{itemize}[leftmargin=1.5em,itemsep=0pt,topsep=2pt]
    \item \texttt{content} y \texttt{vot\_date} opcionales
          (\texttt{minOccurs="0"}), reflejando la realidad del corpus.
    \item \texttt{xs:long} para timestamps Unix (fuera del rango de
          \texttt{xs:integer}).
    \item Patr\'on \texttt{u[0-9]+} para IDs de usuario anonimizados.
\end{itemize}

La coherencia de los resultados XQuery se verific\'o cruzando conteos
obtenidos por distintas v\'ias: el total de preguntas por categor\'ia
(16.952) coincide con el devuelto por \texttt{count(collection(...)//document)}.

%% ============================================
\section{Estado del Desarrollo}
%% ============================================

Todos los componentes se encuentran en estado
\textcolor{statusgreen}{\textbf{Completado}} al 25 de marzo de 2026:
extracci\'on del corpus, XSD original y ampliado, transformaciones XSLT,
carga en eXist-db, cuatro consultas XQuery, dashboard HTML interactivo y
documentaci\'on completa.

%% ============================================
\section{Conclusiones y Trabajo Futuro}
%% ============================================

\subsection{Conclusiones}

\begin{enumerate}[leftmargin=1.8em,itemsep=0pt,topsep=2pt]
    \item El esquema XSD valida los 16.952 documentos documentando
          la opcionalidad real de los campos del corpus.
    \item Las transformaciones XSLT generan visualizaciones HTML
          que permiten explorar el corpus de forma intuitiva.
    \item El an\'alisis de redes confirma la \textbf{baja simetr\'ia}
          (15,4\% participan en ambos roles) y el \textit{clustering}
          elevado de respondedores \textit{top}.
    \item Las comunidades ES e IT son \textbf{pr\'acticamente independientes}
          (solo 2 usuarios en ambas).
    \item La calidad del corpus es alta: 0 inconsistencias de fechas
          y 0 preguntas sin respuesta.
\end{enumerate}

\subsection{Trabajo futuro}

\begin{itemize}[leftmargin=1.5em,itemsep=0pt,topsep=2pt]
    \item \textbf{\'Indices eXist-db}: eliminar \textit{timeouts} en
          consultas de ordenaci\'on global.
    \item \textbf{Anotaci\'on RDF}: enriquecer el XSD con sem\'antica
          para integraci\'on con ontolog\'ias externas.
    \item \textbf{Grafo social}: exportar datos de usuarios a GraphML
          o JSON-LD para visualizaci\'on con Gephi.
    \item \textbf{Comparativa multilingue}: ampliar a franc\'es,
          alem\'an y portugu\'es.
\end{itemize}

%% ============================================
\section{Planificaci\'on temporal}
%% ============================================

\begin{table}[H]
\centering
\caption{Planificaci\'on temporal del proyecto.}
\label{tab:planning}
\begin{tabular}{p{1.8cm}p{2.8cm}l}
\toprule
\textbf{Fase} & \textbf{Actividades} & \textbf{Per\'iodo} \\
\midrule
Preparaci\'on & Repo, eXist-db, XML & 16--17 mar \\
XSD           & Original y ampliado & 17--19 mar \\
Validaci\'on  & validate\_xml.py    & 19 mar \\
XSLT          & ES e IT a HTML      & 19 mar \\
Carga         & cargar\_existdb.py  & 20 mar \\
XQuery        & 4 consultas         & 20--22 mar \\
Dashboard     & ver\_resultados.py  & 24 mar \\
Documentaci\'on & Memoria, LabBook  & 23--25 mar \\
Entrega       & Empaquetado final   & 25 mar \\
\bottomrule
\end{tabular}
\end{table}

%% ============================================
% Referencias
%% ============================================
\begin{thebibliography}{10}

\bibitem{yahoowebscope}
Yahoo! Research,
``Yahoo! Webscope Dataset ydata-yanswers-all-questions-v1\_0,''
Yahoo! Research Alliance Webscope Program, 2007.
[Online]. Disponible: \url{https://webscope.sandbox.yahoo.com/}

\bibitem{yahooanswers-xml}
M. Avil\'es Rodr\'iguez y M. Amagliani,
``YahooAnswers-XML: An\'alisis con XML y eXist-db,''
Repositorio GitHub, 2026.
[Online]. Disponible: \url{https://github.com/matteoamagliani/YahooAnswers-XML}

\bibitem{shen2015}
H. Shen, Z. Li, J. Liu y J.~E. Grant,
``Knowledge Sharing in the Online Social Network of Yahoo! Answers
and Its Implications,''
\textit{IEEE Trans. Comput.}, vol.~64, no.~6, pp. 1715--1728, 2015.

\bibitem{adamic2008}
L.~A. Adamic, J. Zhang, E. Bakshy y M.~S. Ackerman,
``Knowledge Sharing and Yahoo Answers: Everyone Knows Something,''
in \textit{Proc. WWW 2008}, Beijing, China, pp. 665--674.

\bibitem{shah2010}
C. Shah y J. Pomerantz,
``Evaluating and Predicting Answer Quality in Community QA,''
in \textit{Proc. ACM SIGIR 2010}, Geneva, pp. 411--418.

\bibitem{fichman2011}
P. Fichman,
``A Comparative Assessment of Answer Quality on Four Question
Answering Sites,''
\textit{J. Inf. Sci.}, vol.~37, no.~5, pp. 476--486, 2011.

\bibitem{existdb}
eXist Solutions,
``eXist-db Open Source Native XML Database,''
2024. [Online]. Disponible: \url{https://exist-db.org/}

\bibitem{xquery31}
W3C,
``XQuery 3.1: An XML Query Language,''
W3C Recommendation, 2017.
[Online]. Disponible: \url{https://www.w3.org/TR/xquery-31/}

\bibitem{xsd11}
W3C,
``W3C XML Schema Definition Language (XSD) 1.1,''
W3C Recommendation, 2012.
[Online]. Disponible: \url{https://www.w3.org/TR/xmlschema11-1/}

\bibitem{xslt10}
W3C,
``XSL Transformations (XSLT) 1.0,''
W3C Recommendation, 1999.
[Online]. Disponible: \url{https://www.w3.org/TR/xslt}

\end{thebibliography}

\end{document}
